\chapter{Jobs and steps}

Jobs are the unit of parallelism.
Steps are the unit of execution inside a job.

\section{Job dependencies}
Jobs run in parallel by default.
Use \code{needs} to enforce ordering and pass outputs.

\begin{lstlisting}[language=YAML]
jobs:
  build:
    runs-on: ubuntu-latest
    steps:
      - run: ./scripts/build

  test:
    runs-on: ubuntu-latest
    needs: build
    steps:
      - run: ./scripts/test
\end{lstlisting}

\section{Job outputs}
Use outputs to pass data across jobs without recomputation.

\begin{lstlisting}[language=YAML]
jobs:
  meta:
    runs-on: ubuntu-latest
    outputs:
      image_tag: ${{ steps.tag.outputs.value }}
    steps:
      - id: tag
        run: echo "value=sha-${GITHUB_SHA}" >> "$GITHUB_OUTPUT"

  deploy:
    runs-on: ubuntu-latest
    needs: meta
    steps:
      - run: echo "Deploying ${{ needs.meta.outputs.image_tag }}"
\end{lstlisting}

\section{Steps and shells}
Each \code{run} step uses a shell.
On Linux and macOS, it is \code{bash} by default.
On Windows, it is PowerShell.
Declare a different shell if you need one.

\begin{lstlisting}[language=YAML]
- name: Use bash on Windows
  shell: bash
  run: ./scripts/build
\end{lstlisting}

\section{Working directory and environment}
Use \code{working-directory} and \code{env} to keep commands clean.

\begin{lstlisting}[language=YAML]
- name: Run backend tests
  working-directory: services/api
  env:
    NODE_ENV: test
  run: npm test
\end{lstlisting}

\section{Services}
Service containers are useful for databases or queues.
They run alongside your job.

\begin{lstlisting}[language=YAML]
jobs:
  test:
    runs-on: ubuntu-latest
    services:
      postgres:
        image: postgres:15
        env:
          POSTGRES_PASSWORD: postgres
        ports:
          - 5432:5432
    steps:
      - run: ./scripts/test
\end{lstlisting}

\section{Timeouts and resilience}
Use timeouts to avoid hanging workflows.
Use \code{continue-on-error} only for experimental steps.

\begin{lstlisting}[language=YAML]
- name: Optional lints
  continue-on-error: true
  run: ./scripts/lint-extra
\end{lstlisting}

\section{Job design heuristics}
Good jobs are short and explicit:
\begin{itemize}[nosep]
  \item One purpose per job.
  \item One toolchain per job.
  \item Prefer artifacts over shared state.
\end{itemize}

\section{When to split a job}
Split a job when you need isolation, different permissions, or different runners.
Keep steps together when they share a working directory and toolchain.

\section{Retries and flakiness}
If a step is flaky, fix the root cause first.
If you must retry, do it in a script so the workflow remains simple and explicit.

\begin{checklistbox}{Checklist: jobs and steps}
\begin{itemize}[nosep]
  \item Keep jobs focused and composable.
  \item Use \code{needs} instead of ad-hoc ordering.
  \item Prefer outputs over recomputing values.
  \item Set timeouts for long-running jobs.
\end{itemize}
\end{checklistbox}
